% ==============================================================================
% Unified INFORMS LaTeX Master Template (informs4.cls + Online Appendix)
% Supported journal options:
%   opre  - Operations Research           (dblanonrev, \DoubleSpacedXI, <=30pp excl refs)
%   mnsc  - Management Science            (dblanonrev, \OneAndAHalfSpacedXI, <=32pp)
%   msom  - M&SOM                         (dblanonrev, \DoubleSpacedXI, <=32pp all-incl)
%   ijoc  - INFORMS Journal on Computing  (nonblindrev, \DoubleSpacedXI, <=25pp)
%
% NOTE: Download informs4.cls and informs2014.bst from INFORMS Author Portal
% (https://pubsonline.informs.org/authorportal/latex-style-files) into this folder.
% ==============================================================================

\documentclass[mnsc,dblanonrev]{informs4}
% Switch between journals by changing the first option above:
% \documentclass[opre,dblanonrev]{informs4}
% \documentclass[msom,dblanonrev]{informs4}
% \documentclass[ijoc,nonblindrev]{informs4}

% --- Line Spacing (select per target journal) ---
\OneAndAHalfSpacedXI   % Standard for MS (<=32pp)
% \DoubleSpacedXI      % Required for OR, M&SOM (<=32pp all-inclusive), and IJOC

% --- Citations: INFORMS author-year style, no comma between author & year ---
\usepackage{natbib}
\bibpunct[, ]{(}{)}{,}{a}{}{,}
\def\bibfont{\small}
\def\bibsep{\smallskipamount}
\def\bibhang{24pt}
\def\newblock{\ }
\def\BIBand{and}

% --- Core Math, Tables, Figures, Algorithms & Cross-References ---
\usepackage{amsmath,amssymb,amsfonts,mathtools}
\usepackage{booktabs,threeparttable,multirow}
\usepackage{graphicx}
\usepackage{algorithm}
\usepackage{algpseudocode}
\usepackage{xcolor}
% Note for msom: bottom-of-page \footnote is prohibited; use parentheses or endnotes:
% \usepackage{endnotes}
% \let\footnote=\endnote

\usepackage[colorlinks=true,linkcolor=blue!60!black,citecolor=blue!60!black,urlcolor=blue!60!black]{hyperref}

% --- Theorem Environments (numbered sequentially in main text & EC) ---
\TheoremsNumberedThrough
\ECRepeatTheorems
\EquationsNumberedThrough

% --- Clean Macros (Keep minimal; zero orphan macros) ---
\newcommand{\E}{\mathbb{E}}
\newcommand{\R}{\mathbb{R}}

\begin{document}

\RUNAUTHOR{Author Surname} % Hidden automatically under dblanonrev
\RUNTITLE{Short Running Title}

\TITLE{Full Paper Title in Title Case}

\ARTICLEAUTHORS{%
\AUTHOR{First Author}
\AFF{Department, University, City, State ZIP, \EMAIL{author@university.edu}}
\AUTHOR{Second Author}
\AFF{Department, University, City, State ZIP, \EMAIL{coauthor@university.edu}}
}

% --- Abstract ---
% OR: <=200 words, 1 paragraph, no math, no citations.
% MS: <=250 words, 1 paragraph, no math, no citations.
% IJOC: <=300 words, 1 paragraph, no math, no citations.
% M&SOM: <=300 words, 1 paragraph with the 3 bold inline labels below:
%   \textbf{Problem definition:} ...
%   \textbf{Methodology/results:} ...
%   \textbf{Managerial implications:} ...
\ABSTRACT{%
  State the operational or methodological tension, the central model or empirical
  object, the headline mechanism and result magnitude in plain language, and
  the primary regime or managerial implication without equations or citations.
}

\KEYWORDS{keyword one, keyword two, keyword three, keyword four}

% --- Required ONLY for Operations Research (opre) ---
% \SUBJECTCLASS{inventory/production: stochastic; dynamic programming: Markov}
% \AREAOFREVIEW{Operations and Supply Chains}

\maketitle

% ==============================================================================
\section{Introduction}
\label{sec:intro}
% ==============================================================================

% Keep Introduction near-zero notation. Anchor in concrete operational/technical
% friction -> literature positioning -> central object & spine results -> roadmap (if needed).

% ==============================================================================
\section{Model}
\label{sec:model}
% ==============================================================================

% Introduce notation Just-in-Time in dependency order.
% Onboard any specialized concept via: Why -> Intuition -> Definition -> Canonical Term.

% Example Theorem + Proof Sketch + Unified Online Appendix Pointer:
% \begin{proposition}[Threshold Optimality]
% \label{prop:threshold}
% Under stationary arrivals, the optimal policy takes a threshold form.
% \end{proposition}
%
% Prose interpreting the proposition and regime boundary...
% For OR/IJOC (and complex MS/MSOM results), give a 3-6 sentence proof sketch
% naming the bottleneck and load-bearing trick, then point to Online Appendix EC.1:
% The complete proof is provided in Online Appendix~\ref{ec:proofs}.

% ==============================================================================
\section{Numerical Study}
\label{sec:numerics}
% ==============================================================================

% Summarize patterns, mechanisms, and clean rounded magnitudes in prose;
% keep exact decimals inside booktabs tables.

% ==============================================================================
\section{Conclusion}
\label{sec:conclusion}
% ==============================================================================

% Synthesize core takeaway and frame model boundaries positively as future extensions.

% --- Endnotes (Uncomment ONLY for M&SOM if using endnotes instead of parentheses) ---
% \theendnotes

\ACKNOWLEDGMENT{%
  The authors thank the departmental editor, the associate editor, and the
  anonymous referees for their constructive comments.
}

\bibliographystyle{informs2014}
\bibliography{refs}

% ==============================================================================
% UNIFIED ONLINE APPENDIX (100% of proofs, lemmas, tables, extra numerics)
% Note: Page count of Main Text + References ends BEFORE \ECSwitch.
% For M&SOM, Online Appendix must not exceed 16 pages.
% ==============================================================================
\clearpage
\ECSwitch
\ECHead{Online Appendix for ``Full Paper Title in Title Case''}

\section{Proofs of Main Results}
\label{ec:proofs}

\begin{proof}{Proof of Proposition~\ref{prop:threshold}.}
  Write the complete rigorous verification here.
  \Halmos
\end{proof}

\section{Additional Numerical Experiments and Robustness Checks}
\label{ec:numerics}

\end{document}
